\documentclass[../../../include/open-logic-section]{subfiles}

\begin{document}

\olfileid{sfr}{infinite}{card-sb}

\olsection{परिशिष्ट: श्रेडर--बर्नस्टाइनची सिद्धता}

अनौपचारिक संचसिद्धान्ताचा निरोप घेण्यापूर्वी, आणखी एका अनौपचारिक
(पण प्रगत!) सिद्धतेचा विचार करायचा आहे. ही श्रेडर--बर्नस्टाइनची सिद्धता
आहे (\olref[sfr][siz][sb]{thm:schroder-bernstein}): जर
$\cardle{A}{B}$ आणि $\cardle{B}{A}$, तर $\cardeq{A}{B}$; म्हणजे,
!!{injection}s $f \colon A \to B$ आणि $g \colon B \to A$ दिली असता
!!a{bijection} $h \colon A \to B$ असते.

या प्रकरणात आपण डेडेकिंड यांच्या \emph{संवरणांच्या} संकल्पनेचा वापर
केला. किंबहुना याच संकल्पनेचा उपयोग करून डेडेकिंड यांनी
श्रेडर--बर्नस्टाइनची एक सुंदर सिद्धता दिली; ती येथे मांडू. ही सिद्धता
\citet[pp.~157--8]{Potter2004} यांच्या विवेचनाला अगदी जवळून अनुसरते;
तेथील मांडणी किंचित वेगळी पण मूलतः समान आहे. इंटरनेटवर थोडा शोध
घेतल्यास असेही दिसेल की---$\sqrt{2}$ च्या अपरिमेयतेप्रमाणेच---या
प्रमेयाच्या \emph{अनेक} रोचक आणि वेगवेगळ्या सिद्धता आहेत.

\olref[sfr][infinite][dedekind]{Closure} मधील संकेतनासारखेच संकेतन
वापरून, प्रत्येक संच $B$ आणि फलन~$f$ यांसाठी
\[
\Closureofunder{f}{B} = \bigcap \Setabs{X}{B \subseteq X
\text{ आणि $X$ हा $f$-बंद आहे}}
\]
असे मानू. अशा रीतीने परिभाषित केलेला $\Closureofunder{f}{B}$ हा~$B$
चा समावेश असलेला लघुतम $f$-बंद संच आहे; म्हणजे:

\begin{lem}\ollabel{Closureprops}
	कोणतेही फलन $f$ आणि कोणताही $B$ यांसाठी:
	\begin{enumerate}
		\item\ollabel{Closurehaselem} $B \subseteq \Closureofunder{f}{B}$; आणि
		\item\ollabel{Closureclosed} $\Closureofunder{f}{B}$ हा $f$-बंद आहे; आणि
		\item\ollabel{Closuresmallest} $X$ हा $f$-बंद असून $B
		\subseteq X$ असेल, तर $\Closureofunder{f}{B} \subseteq X$.
	\end{enumerate}
\end{lem}

\begin{proof}
\olref[sfr][infinite][dedekind]{closureproperties} मधील सिद्धतेप्रमाणेच.
\end{proof}

श्रेडर--बर्नस्टाइनपर्यंत पोहोचण्यासाठी आपल्याला आणखी एका तथ्याची गरज आहे:

\begin{prop}\ollabel{sbhelper}
जर $A \subseteq B \subseteq C$ आणि $A \approx C$, तर $\cardeq{B}{C}$.
\end{prop}
%READERNOTE{MRINF-003: गोठवलेल्या इंग्रजी निष्कर्षात तुल्यबलता-उक्तीच पहिला संच म्हणून बसवली आहे. गृहीतके आणि पुढील उपयोग यांवरून B आणि C तुल्यबल आहेत हाच अनिवार्य निष्कर्ष मराठीत दिला आहे.}

\begin{proof}
!!a{bijection}  $f \colon C \to A$ दिले असता, $F =
\Closureofunder{f}{C \setminus B}$ असे मानू आणि प्रांत $C$ असलेले फलन
$g$ पुढीलप्रमाणे परिभाषित करू:
\[
	g(x) =
	\begin{cases}
			f(x) &\text{जर $x \in F$}\\
			x & \text{अन्यथा}
	\end{cases}
\]
$g$ हे $C \to B$ असे !!a{bijection} आहे, हे आपण दाखवू. त्यावरून
$\comp{f^{-1}}{g} \colon A \to B$ हे !!a{bijection} आहे आणि सिद्धता
पूर्ण होते.

प्रथम आपला दावा असा आहे की $x \in F$ पण $y\notin F$ असेल, तर $g(x) \neq
g(y)$. विरोधाभासाने सिद्ध करण्यासाठी तसे नाही असे समजा; मग $y = g(y) =
g(x) = f(x)$. $x \in F$ आहे आणि \olref{Closureprops} नुसार $F$ हा
$f$-बंद आहे, म्हणून $y = f(x) \in F$; हा विरोधाभास आहे.

आता $g(x) = g(y)$ असे समजा. वरील दाव्यानुसार, $x \in F$ तेव्हा आणि
केवळ तेव्हाच, जेव्हा $y \in F$. जर $x, y \in F$, तर $f(x) = g (x) =
g(y) = f(y)$; आणि $f$ हे !!a{bijection} असल्यामुळे $x = y$. जर
$x, y \notin F$, तर $x = g(x) = g(y) = y$. म्हणून $g$ हे
!!a{injection} आहे.

$\ran{g} = B$ हे दाखवणे उरते. म्हणून $x \in B \subseteq C$ स्थिर करू.
जर $x \notin F$, तर $g(x) = x$. जर $x \in F$, तर कोणत्यातरी
$y \in F$ साठी $x = f(y)$; कारण तसे नसते, तर $F \setminus \{x\}$ हा
$f$-बंद संच असता आणि त्यात $C\setminus B$ चा समावेश असता; पण
\olref{Closureprops} नुसार हे अशक्य आहे. आता $g(y) = f(y) = x$.
\end{proof}

शेवटी, मुख्य निष्कर्षाची सिद्धता पुढीलप्रमाणे. फलन $h$ आणि संच $D$ दिले
असता आपण $\funimage{h}{D} = \Setabs{h(x)}{x \in D}$ अशी व्याख्या करतो,
हे आठवा.

\begin{proof}[श्रेडर--बर्नस्टाइनची सिद्धता] $f \colon A \to B$
आणि $g \colon B \to A$ ही !!{injection}s असू द्या.
$\funimage{f}{A} \subseteq B$ असल्यामुळे
$\funimage{g}{\funimage{f}{A}} \subseteq g[B] \subseteq A$.
तसेच, $\comp{f}{g} \colon A \to
\funimage{g}{\funimage{f}{A}}$ हे !!{injection} आहे, कारण $g$ आणि
$f$ दोन्ही तशी आहेत; आणि त्याचा सहप्रांत ज्या रीतीने परिभाषित केला आहे,
त्यावरून $\comp{f}{g}$ हे खरे तर !!a{bijection} आहे. म्हणून
$\cardeq{\funimage{g}{\funimage{f}{A}}}{A}$; आणि त्यामुळे
\olref{sbhelper} नुसार $h \colon A \to \funimage{g}{B}$ असे
!!a{bijection} आहे. शिवाय, $g^{-1}$ हे
$\funimage{g}{B} \to B$ असे !!a{bijection} आहे. म्हणून
$\comp{h}{g^{-1}} \colon A \to B$ हे !!a{bijection} आहे.
\end{proof}

\end{document}
